\documentclass[12pt]{article}
\usepackage[utf8]{inputenc}
\usepackage[T1]{fontenc}
\usepackage{hyperref}
\title{Amon Language Resources for K-12 Curriculum Development}
\author{Compiled by Claude Sonnet 5 (claude-sonnet-5) for Andr\'as Kornai}
\date{}
\begin{document}
\maketitle
\section{Introduction}
Amon (also called Umon; ISO 639-3 umm) is a language of Nigeria. My search found it described as an Upper Cross River language of Cross River State. This is not the Amo (Map, Timap) language of the Kaduna and Plateau area, which has a different ISO code and its own New Testament. Search turned up almost nothing specific to Amon, so this report is short.
\section{Primary Educational Resources}
I found no primers, textbooks or school materials for Amon.
\section{Linguistic and Reference Materials}
Only an encyclopedic overview appeared: \url{https://wikipedia.com/wiki/ISO_639:umm}. A language-code catalogue entry also appeared: \url{https://finto.fi/lexvo/en/page/umm}.
\section{Religious and Community Resources}
I found no Amon-specific scripture or community organization. The Amo New Testament (\url{https://dbs.org/bibles/AMODWC}) belongs to a different language and should not be used for Amon.
\section{Digital and Audio Resources}
I found no audio, video or corpus resources for Amon.
\section{Curriculum Development Context}
Little is documented online. A curriculum would need new fieldwork, community orthography decisions and partnership with local speakers and education authorities. The reference description says the language is not known to be taught in schools.
\end{document}